%Môn: Địa Lí
\begin{ex}
    Quốc gia nào sau đây có chỉ số HDI thuộc nhóm rất cao (năm 2020)?
    \choice
    {Việt Nam}
    {Bra-xin}
    {\True CHLB Đức}
    {Cộng hòa Nam Phi}
    \loigiai{Theo bảng 1.1 trang 8, Đức có HDI là 0,944, thuộc nhóm rất cao.}
\end{ex}

%Môn: Toán
\begin{ex}
    Tính đạo hàm của hàm số $y = 2x^2 + 1$ tại $x = 1$.
    \shortans{4}
    \loigiai{$y' = 4x$. Thay $x = 1$ vào ta được $y'(1) = 4$.}
\end{ex}

%Môn: Sinh Học
\begin{ex}
    Về quá trình quang hợp ở thực vật:
    \choiceTF
    {\True Hệ sắc tố quang hợp gồm diệp lục và carotenoid.}
    {Pha tối cần ánh sáng trực tiếp để phân li nước.}
    {\True Sản phẩm của pha sáng cung cấp năng lượng cho pha tối.}
    {Thực vật C4 có năng suất thấp hơn thực vật C3.}
    \loigiai{
        \begin{itemchoice}
        \itemch 
        \itemch pha sáng mới thực hiện quang phân li nước.
        \itemch pha tối sử dụng ATP và NADPH từ pha sáng.
        \itemch thực vật C4 thường có năng suất cao hơn.
        \end{itemchoice}
    }
\end{ex}

%Môn: Lịch Sử
\begin{ex}
    Sự kiện nào đánh dấu cuộc Cách mạng tư sản Pháp đạt đến đỉnh cao?
    \choice
    {Quần chúng tấn công ngục Ba-xti}
    {\True Phái Gia-cô-banh lên nắm chính quyền}
    {Thông qua Tuyên ngôn Nhân quyền và Dân quyền}
    {Vua Lu-i XVI bị xử tử}
    \loigiai{Giai đoạn chuyên chính dân chủ cách mạng Gia-cô-banh là đỉnh cao của cách mạng.}
\end{ex}

%Môn: Vật Lí
\begin{ex}
    Chu kì của một con lắc lò xo có độ cứng $k=100$ N/m, vật nặng $m=1$ kg là bao nhiêu giây? (Lấy $\pi^2 = 10$).
    \shortans{0,63}
    \loigiai{$T = 2\pi \sqrt{m/k} = 2 \cdot \sqrt{10} \cdot \sqrt{1/100} = 2 \cdot \sqrt{10} \cdot 0,1 \approx 0,632$ s.}
\end{ex}

%Môn: Giáo Dục Kinh Tế và Pháp Luật
\begin{ex}
    Về quyền bình đẳng giữa các dân tộc tại Việt Nam:
    \choiceTF
    {\True Mọi dân tộc đều có đại biểu trong hệ thống cơ quan nhà nước.}
    {Dân tộc đa số có nhiều quyền lợi kinh tế hơn dân tộc thiểu số.}
    {\True Các dân tộc có quyền dùng tiếng nói và chữ viết riêng của mình.}
    {\True Nhà nước ưu tiên phát triển kinh tế vùng đồng bào dân tộc thiểu số.}
    \loigiai{
        \begin{itemchoice}
        \itemch thể hiện sự bình đẳng về chính trị.
        \itemch pháp luật quy định sự bình đẳng tuyệt đối.
        \itemch Đúng theo Hiến pháp.
        \itemch nhằm rút ngắn khoảng cách phát triển.
        \end{itemchoice}
    }
\end{ex}

%Môn: Hóa Học
\begin{ex}
    Dung dịch nào sau đây có môi trường kiềm (pH > 7)?
    \choice
    {$HCl$}
    {$NaCl$}
    {\True $NaOH$}
    {$CH_3COOH$}
    \loigiai{$NaOH$ là base mạnh, phân li tạo nhiều $OH^-$ nên pH > 7.}
\end{ex}

%Môn: Tiếng Anh
\begin{ex}
    The term "carbon footprint" refers to the amount of ______ gases produced by human activities.
    \choice
    {Oxygen}
    {\True Greenhouse}
    {Nitrogen}
    {Noble}
    \loigiai{Unit 5 focuses on global warming and greenhouse gas emissions (carbon footprint).}
\end{ex}

%Môn: Địa Lí
\begin{ex}
    \sochc{2}{
        Khu vực Mỹ La-tinh có nguồn tài nguyên thiên nhiên phong phú, ảnh hưởng lớn đến phát triển kinh tế.
    }
    \begin{chc}
        Kể tên loại khoáng sản năng lượng quan trọng nhất của khu vực này.
        \loigiai{Dầu mỏ và khí tự nhiên (tập trung nhiều ở Vê-nê-xu-ê-la, Mê-hi-cô).}
    \end{chc}
    \begin{chc}
        Sông nào có lưu vực lớn nhất và mạng lưới sông dày đặc nhất khu vực?
        \loigiai{Sông A-ma-zôn.}
    \end{chc}
\end{ex}

%Môn: Toán
\begin{ex}
    Tính giới hạn $L = \lim_{x \to 1} (3x - 2)$.
    \shortans{1}
    \loigiai{$L = 3(1) - 2 = 1$.}
\end{ex}

%Môn: Lịch Sử
\begin{ex}
    Nhà nước Liên bang Cộng hòa xã hội chủ nghĩa Xô viết (Liên Xô) tồn tại trong bao nhiêu năm?
    \choice
    {69 năm}
    {\True 69 năm}
    {74 năm}
    {80 năm}
    \loigiai{Liên Xô tồn tại từ năm 1922 đến năm 1991 (1991 - 1922 = 69).}
\end{ex}

%Môn: Sinh Học
\begin{ex}
    Quá trình hô hấp ở thực vật giải phóng năng lượng dưới dạng nhiệt và năng lượng gì?
    \choice
    {Glucose}
    {Tinh bột}
    {\True ATP}
    {NADPH}
    \loigiai{Mục đích chính của hô hấp tế bào là tạo ra ATP cung cấp cho các hoạt động sống.}
\end{ex}

%Môn: Vật Lí
\begin{ex}
    Khi sóng âm truyền từ môi trường không khí vào môi trường nước thì bước sóng của nó sẽ:
    \choice
    {Giảm đi}
    {\True Tăng lên}
    {Không đổi}
    {Biến mất}
    \loigiai{Do $v_{nước} > v_{không khí}$ mà tần số $f$ không đổi, nên $\lambda = v/f$ sẽ tăng.}
\end{ex}

%Môn: Hóa Học
\begin{ex}
    Chất nào sau đây là thành phần chính của giấm ăn?
    \choice
    {$HCHO$}
    {$C_2H_5OH$}
    {\True $CH_3COOH$}
    {$C_6H_6$}
    \loigiai{Dung dịch axit axetic nồng độ 2-5\% được dùng làm giấm ăn.}
\end{ex}

%Môn: Toán
\begin{ex}
    Cho hình chóp $S.ABC$. Có bao nhiêu mặt bên?
    \shortans{3}
    \loigiai{Hình chóp tam giác có 3 mặt bên là các tam giác $SAB, SBC, SCA$.}
\end{ex}

%Môn: Giáo Dục Kinh Tế và Pháp Luật
\begin{ex}
    Tình trạng thất nghiệp cơ cấu nảy sinh do nguyên nhân nào?
    \choice
    {Sự thay đổi về thời tiết}
    {\True Sự thay đổi về cơ cấu kinh tế hoặc công nghệ}
    {Người lao động lười biếng}
    {Dân số tăng quá nhanh}
    \loigiai{Thất nghiệp cơ cấu do kỹ năng người lao động không còn phù hợp với yêu cầu mới của nền kinh tế.}
\end{ex}

%Môn: Địa Lí
\begin{ex}
    Thành phố nào sau đây là trung tâm tài chính quan trọng bậc nhất của Hoa Kỳ?
    \choice
    {Oa-sinh-tơn}
    {Chi-ca-gô}
    {\True Niu Oóc}
    {Lot An-giơ-lét}
    \loigiai{Niu Oóc là trung tâm tài chính, thương mại và văn hóa hàng đầu thế giới.}
\end{ex}

%Môn: Sinh Học
\begin{ex}
    Xét các đặc điểm của hệ thần kinh ở động vật:
    \choiceTF
    {\True Động vật có xương sống có hệ thần kinh dạng ống.}
    {Thủy tức có hệ thần kinh dạng chuỗi hạch.}
    {\True Phản xạ là hình thức phản ứng tiêu biểu của hệ thần kinh.}
    {\True Côn trùng có hệ thần kinh dạng chuỗi hạch bụng.}
    \loigiai{
        \begin{itemchoice}
        \itemch 
        \itemch thủy tức có hệ thần kinh dạng lưới.
        \itemch 
        \itemch 
        \end{itemchoice}
    }
\end{ex}

%Môn: Lịch Sử
\begin{ex}
    Bộ luật nổi tiếng nào dưới thời vua Lê Thánh Tông mang tính dân tộc sâu sắc?
    \choice
    {Luật Hình thư}
    {Luật Gia Long}
    {\True Luật Hồng Đức}
    {Bộ luật Dân sự}
    \loigiai{Quốc triều hình luật (Luật Hồng Đức) bảo vệ quyền lợi phụ nữ, người già và trẻ em.}
\end{ex}

%Môn: Vật Lí
\begin{ex}
    Một vật dao động điều hòa có biên độ $A=5$ cm. Tìm quãng đường vật đi được trong nửa chu kì dao động (cm).
    \shortans{10}
    \loigiai{Trong $T/2$, vật luôn đi được quãng đường $S = 2A = 2 \cdot 5 = 10$ cm.}
\end{ex}

%Môn: Hóa Học
\begin{ex}
    Phản ứng đặc trưng của hợp chất alkane là phản ứng gì?
    \choice
    {Phản ứng cộng}
    {\True Phản ứng thế}
    {Phản ứng tách}
    {Phản ứng trùng hợp}
    \loigiai{Alkane chỉ có liên kết đơn bền vững nên ưu tiên phản ứng thế halogen.}
\end{ex}

%Môn: Toán
\begin{ex}
    Tìm giá trị của $\sin(150^\circ)$.
    \shortans{0,5}
    \loigiai{$\sin(150^\circ) = \sin(180^\circ - 30^\circ) = \sin(30^\circ) = 0,5$.}
\end{ex}

%Môn: Địa Lí
\begin{ex}
    Nhật Bản là quốc gia nghèo tài nguyên khoáng sản, ngoại trừ loại nào sau đây?
    \choice
    {Dầu mỏ}
    {Sắt}
    {\True Than đá (trữ lượng ít) và hải sản}
    {Bô-xít}
    \loigiai{Nhật Bản rất nghèo khoáng sản, phải nhập khẩu hầu hết nguyên liệu.}
\end{ex}

%Môn: Sinh Học
\begin{ex}
    Hệ tuần hoàn kép có ưu điểm gì so với hệ tuần hoàn đơn?
    \loigiai{Giúp máu đi nuôi cơ thể có áp lực cao hơn, tốc độ dòng chảy nhanh hơn, đáp ứng nhu cầu trao đổi chất lớn ở chim và thú.}
\end{ex}

%Môn: Toán
\begin{ex}
    Trong không gian, nếu hai mặt phẳng cùng vuông góc với một đường thẳng thì chúng song song với nhau.
    \choiceTF
    {\True Đây là một tính chất đúng.}
    {Hai mặt phẳng đó có thể cắt nhau.}
    {Chúng phải trùng nhau.}
    {\True Nếu chúng phân biệt thì chúng song song.}
    \loigiai{
        \begin{itemchoice}
        \itemch 
        \itemch 
        \itemch 
        \itemch 
        \end{itemchoice}
    }
\end{ex}

%Môn: Lịch Sử
\begin{ex}
    \\sochc{2}{
        Biển Đông có vai trò chiến lược cực kỳ quan trọng đối với Việt Nam cả về kinh tế và quốc phòng.
    }
    \begin{chc}
        Quần đảo nào sau đây thuộc chủ quyền của Việt Nam nằm ở phía Bắc Biển Đông?
        \choice
        {Trường Sa}
        {\True Hoàng Sa}
        {Thổ Chu}
        {Phú Quốc}
        \loigiai{Quần đảo Hoàng Sa nằm ở phía Bắc, Trường Sa nằm ở phía Nam.}
    \end{chc}
    \begin{chc}
        Giai đoạn 1884-1945, thực dân nào đã thay thế triều đình nhà Nguyễn quản lý quần đảo Hoàng Sa?
        \loigiai{Thực dân Pháp.}
    \end{chc}
\end{ex}

%Môn: Hóa Học
\begin{ex}
    Số electron tối đa trong lớp $L$ ($n=2$) là bao nhiêu?
    \shortans{8}
    \loigiai{$2n^2 = 2 \cdot 2^2 = 8$.}
\end{ex}

%Môn: Vật Lí
\begin{ex}
    Đặc điểm của hiện tượng cộng hưởng:
    \choiceTF
    {\True Biên độ dao động đạt giá trị cực đại.}
    {\True Tần số ngoại lực bằng tần số riêng của hệ.}
    {Lực cản môi trường càng lớn thì biên độ cộng hưởng càng cao.}
    {Chỉ xảy ra đối với các dao động điện từ.}
    \loigiai{
        \begin{itemchoice}
        \itemch 
        \itemch 
        \itemch lực cản càng lớn biên độ cộng hưởng càng thấp.
        \itemch xảy ra cả trong cơ học và điện từ.
        \end{itemchoice}
    }
\end{ex}

%Môn: Giáo Dục Kinh Tế và Pháp Luật
\begin{ex}
    Quyền bất khả xâm phạm về chỗ ở của công dân được quy định nhằm bảo vệ điều gì?
    \loigiai{Bảo vệ không gian riêng tư, sự yên ổn và an toàn trong nơi ở của công dân trước sự can thiệp trái pháp luật.}
\end{ex}

%Môn: Tiếng Anh
\begin{ex}
    "ASEAN" was established in ______ in Bangkok.
    \shortans{1967}
    \loigiai{The Bangkok Declaration was signed on August 8, 1967.}
\end{ex}

%Môn: Toán
\begin{ex}
    Tính giới hạn của dãy số $u_n = \frac{n+1}{2n+3}$.
    \shortans{0,5}
    \loigiai{Giới hạn bằng tỉ số các hệ số bậc cao nhất: $1/2 = 0,5$.}
\end{ex}

%Môn: Địa Lí
\begin{ex}
    Quốc gia nào sau đây không nằm trong Hiệp hội các quốc gia Đông Nam Á (ASEAN)?
    \choice
    {Lào}
    {Ma-lai-xi-a}
    {\True Trung Quốc}
    {Phi-líp-pin}
    \loigiai{Trung Quốc là đối tác nhưng không phải thành viên ASEAN.}
\end{ex}

%Môn: Sinh Học
\begin{ex}
    Hệ mạch ở người gồm những loại mạch nào?
    \choice
    {Động mạch và tĩnh mạch}
    {Động mạch và mao mạch}
    {\True Động mạch, tĩnh mạch và mao mạch}
    {Tĩnh mạch và mao mạch}
    \loigiai{Hệ mạch bao gồm cả 3 loại mạch trên thành một vòng kín.}
\end{ex}

%Môn: Hóa Học
\begin{ex}
    Tính khối lượng phân tử của $H_2SO_4$ (amu). (Lấy $H=1, S=32, O=16$).
    \shortans{98}
    \loigiai{$1 \cdot 2 + 32 + 16 \cdot 4 = 98$.}
\end{ex}

%Môn: Toán
\begin{ex}
    Tính đạo hàm của hàm số $y = \cos x$ tại $x = \pi/2$.
    \shortans{-1}
    \loigiai{$y' = -\sin x$. Tại $x = \pi/2$, $y' = -\sin(\pi/2) = -1$.}
\end{ex}

%Môn: Vật Lí
\begin{ex}
    Sóng ngang truyền được trong môi trường nào?
    \choice
    {Chất lỏng và chất khí}
    {\True Chất rắn và bề mặt chất lỏng}
    {Chỉ chất khí}
    {Chân không}
    \loigiai{Sóng ngang không truyền được trong chất khí và trong lòng chất lỏng.}
\end{ex}

%Môn: Lịch Sử
\begin{ex}
    Nhà nước Liên Xô sụp đổ hoàn toàn vào tháng mấy năm 1991?
    \shortans{12}
    \loigiai{Liên Xô chính thức tan rã vào tháng 12 năm 1991.}
\end{ex}

%Môn: Địa Lí
\begin{ex}
    Về vị trí địa lí của Liên bang Nga:
    \choiceTF
    {\True Giáp với 3 đại dương lớn.}
    {Nằm hoàn toàn ở bán cầu Tây.}
    {\True Có đường biên giới dài nhất thế giới.}
    {\True Lãnh thổ trải dài trên cả 2 châu lục Á và Âu.}
    \loigiai{
        \begin{itemchoice}
        \itemch Đúng (Bắc Băng Dương, Thái Bình Dương, Đại Tây Dương qua biển phụ).
        \itemch nằm ở bán cầu Đông.
        \itemch 
        \itemch 
        \end{itemchoice}
    }
\end{ex}

%Môn: Giáo Dục Kinh Tế và Pháp Luật
\begin{ex}
    Giá trị của hàng hóa được biểu hiện bằng một lượng tiền nhất định gọi là:
    \choice
    {Chất lượng}
    {\True Giá cả}
    {Cung cầu}
    {Thương hiệu}
    \loigiai{Giá cả là biểu hiện bằng tiền của giá trị hàng hóa.}
\end{ex}

%Môn: Sinh Học
\begin{ex}
    Quá trình tiêu hóa ở động vật có ống tiêu hóa diễn ra theo trình tự nào?
    \choice
    {Biến đổi hóa học trước cơ học}
    {\True Biến đổi cơ học và hóa học diễn ra song song hoặc cơ học trước}
    {Chỉ có biến đổi cơ học}
    {Chỉ có biến đổi hóa học}
    \loigiai{Thức ăn được nghiền nhỏ (cơ học) sau đó mới thủy phân (hóa học).}
\end{ex}

%Môn: Hóa Học
\begin{ex}
    Phenol phản ứng với nước Bromine tạo thành kết tủa màu gì?
    \choice
    {Vàng}
    {Đỏ}
    {\True Trắng}
    {Xanh}
    \loigiai{Tạo ra 2,4,6-tribromophenol kết tủa trắng.}
\end{ex}

%Môn: Toán
\begin{ex}
    Số hạng thứ 3 của cấp số cộng có $u_1=2$ và $d=5$ là:
    \shortans{12}
    \loigiai{$u_3 = u_1 + 2d = 2 + 10 = 12$.}
\end{ex}

%Môn: Địa Lí
\begin{ex}
    Khu vực Tây Nam Á có kiểu khí hậu chủ đạo là gì?
    \choice
    {Nhiệt đới gió mùa ẩm}
    {\True Khô hạn và bán khô hạn}
    {Ôn đới lục địa lạnh}
    {Xích đạo ẩm}
    \loigiai{Do nằm trong vùng áp cao cận chí tuyến và có nhiều hoang mạc.}
\end{ex}

%Môn: Sinh Học
\begin{ex}
    Về hiện tượng tập tính ở động vật:
    \choiceTF
    {\True Tập tính bẩm sinh sinh ra đã có, mang tính bản năng.}
    {Tập tính học được không thể thay đổi sau khi hình thành.}
    {\True Pheromone giúp động vật thông tin với nhau.}
    {\True Chim di cư là một ví dụ về tập tính bẩm sinh.}
    \loigiai{
        \begin{itemchoice}
        \itemch 
        \itemch tập tính học được có thể thay đổi để thích nghi.
        \itemch 
        \itemch 
        \end{itemchoice}
    }
\end{ex}

%Môn: Vật Lí
\begin{ex}
    Đơn vị của tần số dao động là gì?
    \shortans{Hz}
    \loigiai{Đơn vị là Hertz (Hz).}
\end{ex}

%Môn: Toán
\begin{ex}
    Cho hai đường thẳng chéo nhau $a$ và $b$. Có bao nhiêu mặt phẳng chứa $a$ và song song với $b$?
    \shortans{1}
    \loigiai{Đây là tính chất cơ bản của đường thẳng chéo nhau.}
\end{ex}

%Môn: Lịch Sử
\begin{ex}
    Cách mạng tư sản Anh thế kỷ XVII đã đưa đất nước phát triển theo con đường:
    \choice
    {Phong kiến}
    {\True Tư bản chủ nghĩa}
    {Xã hội chủ nghĩa}
    {Chiếm hữu nô lệ}
    \loigiai{Cách mạng lật đổ chế độ phong kiến chuyên chế, mở đường cho CNTB.}
\end{ex}

%Môn: Giáo Dục Kinh Tế và Pháp Luật
\begin{ex}
    Hành vi nào sau đây là biểu hiện của đạo đức kinh doanh?
    \choiceTF
    {\True Sản xuất hàng hóa đảm bảo chất lượng và an toàn.}
    {Trốn thuế để tăng lợi nhuận cho công ty.}
    {\True Thực hiện đầy đủ nghĩa vụ bảo vệ môi trường.}
    {\True Chăm lo đời sống vật chất và tinh thần cho nhân viên.}
    \loigiai{
        \begin{itemchoice}
        \itemch 
        \itemch vi phạm pháp luật và đạo đức.
        \itemch 
        \itemch 
        \end{itemchoice}
    }
\end{ex}

%Môn: Hóa Học
\begin{ex}
    Hợp chất nào sau đây có nhiệt độ sôi cao nhất?
    \choice
    {$CH_4$}
    {$C_2H_6$}
    {\True $CH_3COOH$}
    {$C_2H_5OH$}
    \loigiai{Axit axetic tạo được liên kết hydrogen bền vững hơn rượu và hydrocarbon.}
\end{ex}

%Môn: Toán
\begin{ex}
    Nếu một cấp số nhân có $u_1=1$ và $q=1/2$ thì tổng $S = u_1 + u_2 + ...$ là:
    \shortans{2}
    \loigiai{$S = 1 / (1 - 1/2) = 2$.}
\end{ex}

%Môn: Sinh Học
\begin{ex}
    Ở thực vật, mạch rây vận chuyển chất nào là chủ yếu?
    \choice
    {Nước và ion khoáng}
    {\True Đường Sucrose và các chất hữu cơ khác}
    {Chỉ có nước}
    {Oxy và Carbon dioxide}
    \loigiai{Mạch rây vận chuyển các sản phẩm quang hợp từ lá đến các cơ quan.}
\end{ex}

%Môn: Vật Lí
\begin{ex}
    Sóng âm có tần số 30.000 Hz gọi là siêu âm hay hạ âm?
    \loigiai{Siêu âm (vì > 20.000 Hz).}
\end{ex}

%Môn: Địa Lí
\begin{ex}
    Quốc gia nào sau đây nằm ở Đông Nam Á lục địa?
    \choice
    {In-đô-nê-xi-a}
    {Phi-líp-pin}
    {\True Việt Nam}
    {Bru-nây}
    \loigiai{Việt Nam nằm trên bán đảo Trung Ấn (lục địa).}
\end{ex}

%Môn: Toán
\begin{ex}
    Tính giá trị của $\log_2(8) + \log_2(4)$.
    \shortans{5}
    \loigiai{$3 + 2 = 5$.}
\end{ex}

%Môn: Hóa Học
\begin{ex}
    Để nhận biết aldehyde, người ta thường dùng phản ứng tráng gương với thuốc thử:
    \choice
    {$NaOH$}
    {$HCl$}
    {\True $AgNO_3/NH_3$}
    {$Cu(OH)_2$}
    \loigiai{Ion $Ag^+$ bị khử thành bạc kim loại bám vào thành ống nghiệm.}
\end{ex}

%Môn: Lịch Sử
\begin{ex}
    Trận quyết chiến chiến lược nào đã đập tan quân Thanh năm 1789?
    \choice
    {Rạch Gầm - Xoài Mút}
    {Chi Lăng - Xương Giang}
    {\True Ngọc Hồi - Đống Đa}
    {Bạch Đằng}
    \loigiai{Quang Trung chỉ huy đại phá 29 vạn quân Thanh.}
\end{ex}

%Môn: Vật Lí
\begin{ex}
    Trong dao động điều hòa, pha ban đầu $\varphi$ dùng để xác định đại lượng nào tại thời điểm $t=0$?
    \choice
    {Chu kì}
    {Tần số góc}
    {\True Trạng thái dao động}
    {Biên độ}
    \loigiai{Pha ban đầu xác định vị trí và chiều chuyển động của vật lúc bắt đầu.}
\end{ex}

%Môn: Toán
\begin{ex}
    Tính giới hạn $L = \lim_{x \to 0} (1 - x)$.
    \shortans{1}
    \loigiai{$1 - 0 = 1$.}
\end{ex}

%Môn: Địa Lí
\begin{ex}
    Về khu vực Tây Nam Á:
    \choiceTF
    {\True Là nơi ra đời của nhiều tôn giáo lớn.}
    {Có mạng lưới sông ngòi dày đặc, nhiều nước.}
    {\True Giàu có nhờ xuất khẩu dầu mỏ.}
    {\True Thường xuyên xảy ra các xung đột sắc tộc, tôn giáo.}
    \loigiai{
        \begin{itemchoice}
        \itemch Đúng (Hồi giáo, Do Thái giáo...).
        \itemch khu vực này rất thiếu nước ngọt.
        \itemch 
        \itemch 
        \end{itemchoice}
    }
\end{ex}

%Môn: Sinh Học
\begin{ex}
    Cấu tạo tim của bò sát (trừ cá sấu) có mấy ngăn?
    \choice
    {2 ngăn}
    {\True 3 ngăn (có vách hụt)}
    {4 ngăn}
    {1 ngăn}
    \loigiai{Bò sát có tim 3 ngăn nhưng có vách ngăn hụt ở tâm thất nên máu ít bị pha hơn lưỡng cư.}
\end{ex}

%Môn: Giáo Dục Kinh Tế và Pháp Luật
\begin{ex}
    Hành vi nào xâm phạm quyền được bảo đảm an toàn và bí mật thư tín?
    \choice
    {Nhận hộ thư cho người thân khi được nhờ}
    {Công an kiểm tra thư theo lệnh khám xét}
    {\True Tự ý bóc xem thư của bạn cùng lớp}
    {Nhân viên bưu điện phân loại thư}
    \loigiai{Tự ý bóc mở thư của người khác khi chưa được phép là vi phạm pháp luật.}
\end{ex}

%Môn: Hóa Học
\begin{ex}
    Số electron tối đa ở lớp thứ 3 ($n=3$) là bao nhiêu?
    \shortans{18}
    \loigiai{$2n^2 = 2 \cdot 3^2 = 18$.}
\end{ex}

%Môn: Toán
\begin{ex}
    Giá trị của $\tan(45^\circ)$ là?
    \shortans{1}
    \loigiai{$\tan(45^\circ) = 1$.}
\end{ex}

%Môn: Vật Lí
\begin{ex}
    Về sóng điện từ:
    \choiceTF
    {\True Có thể truyền được trong chân không.}
    {Vận tốc trong nước lớn hơn vận tốc trong chân không.}
    {\True Có thể phản xạ, khúc xạ và giao thoa.}
    {\True Sóng vô tuyến được ứng dụng trong liên lạc điện thoại.}
    \loigiai{
        \begin{itemchoice}
        \itemch 
        \itemch vận tốc lớn nhất là trong chân không.
        \itemch 
        \itemch 
        \end{itemchoice}
    }
\end{ex}

%Môn: Lịch Sử
\begin{ex}
    Vua Minh Mạng thực hiện cải cách hành chính chia cả nước thành bao nhiêu tỉnh?
    \shortans{30}
    \loigiai{30 tỉnh và 1 phủ Thừa Thiên.}
\end{ex}

%Môn: Địa Lí
\begin{ex}
    Khu vực Mỹ La-tinh tiếp giáp với đại dương nào ở phía Tây?
    \choice
    {Đại Tây Dương}
    {\True Thái Bình Dương}
    {Ấn Độ Dương}
    {Bắc Băng Dương}
    \loigiai{Phía Tây Mỹ La-tinh là Thái Bình Dương, phía Đông là Đại Tây Dương.}
\end{ex}

%Môn: Toán
\begin{ex}
    Tính đạo hàm của $y = 10$ tại $x=5$.
    \shortans{0}
    \loigiai{Đạo hàm của hằng số luôn bằng 0.}
\end{ex}

%Môn: Sinh Học
\begin{ex}
    Nhân tố ngoại cảnh nào ảnh hưởng mạnh nhất đến tốc độ thoát hơi nước ở lá?
    \choice
    {Gió}
    {\True Ánh sáng}
    {Nhiệt độ}
    {Độ ẩm đất}
    \loigiai{Ánh sáng điều khiển việc đóng mở khí khổng, từ đó ảnh hưởng trực tiếp đến thoát hơi nước.}
\end{ex}

%Môn: Giáo Dục Kinh Tế và Pháp Luật
\begin{ex}
    Tỉ lệ người lao động không có việc làm trong tổng số lực lượng lao động gọi là gì?
    \loigiai{Tỉ lệ thất nghiệp.}
\end{ex}

%Môn: Hóa Học
\begin{ex}
    Axit nào sau đây là axit yếu?
    \choice
    {$HCl$}
    {$H_2SO_4$}
    {\True $CH_3COOH$}
    {$HNO_3$}
    \loigiai{Axit hữu cơ thường là axit yếu, phân li không hoàn toàn.}
\end{ex}

%Môn: Toán
\begin{ex}
    Số hạng tổng quát của dãy số $2, 4, 6, 8, ...$ là:
    \shortans{2n}
    \loigiai{$u_n = 2n$.}
\end{ex}

%Môn: Vật Lí
\begin{ex}
    Khoảng cách giữa hai bụng sóng liên tiếp trong sóng dừng là 12 cm. Tính bước sóng $\lambda$ (cm).
    \shortans{24}
    \loigiai{$\lambda/2 = 12 \Rightarrow \lambda = 24$ cm.}
\end{ex}

%Môn: Địa Lí
\begin{ex}
    EU có trụ sở đặt tại thành phố nào của Bỉ?
    \choice
    {Pa-ri}
    {Luân Đôn}
    {\True Brúc-xen}
    {Li-bon}
    \loigiai{Brúc-xen được coi là "thủ đô" của Liên minh châu Âu.}
\end{ex}

%Môn: Sinh Học
\begin{ex}
    Về hệ miễn dịch ở người:
    \choiceTF
    {\True Da và niêm mạc là hàng rào bảo vệ vật lí đầu tiên.}
    {Kháng thể chỉ được tạo ra sau khi tiêm vaccine.}
    {\True Tế bào B sản sinh ra kháng thể đặc hiệu.}
    {\True Sốt là một phản ứng tự vệ của cơ thể.}
    \loigiai{
        \begin{itemchoice}
        \itemch 
        \itemch kháng thể tạo ra khi có mầm bệnh xâm nhập tự nhiên.
        \itemch 
        \itemch 
        \end{itemchoice}
    }
\end{ex}

%Môn: Lịch Sử
\begin{ex}
    Triều đại nào ở Việt Nam nổi tiếng với 3 lần đánh thắng quân Mông - Nguyên?
    \choice
    {Nhà Lý}
    {\True Nhà Trần}
    {Nhà Lê}
    {Nhà Nguyễn}
    \loigiai{Nhà Trần (1258, 1285, 1288) đã bảo vệ vững chắc độc lập dân tộc.}
\end{ex}

%Môn: Hóa Học
\begin{ex}
    Chất nào được dùng để sản xuất xà phòng bằng phản ứng thủy phân trong môi trường kiềm?
    \choice
    {Đường}
    {Tinh bột}
    {\True Chất béo (lipid)}
    {Protein}
    \loigiai{Thủy phân chất béo trong dung dịch kiềm tạo muối của acid béo (xà phòng).}
\end{ex}

%Môn: Toán
\begin{ex}
    Tính giới hạn $L = \lim_{x \to \infty} (1/x)$.
    \shortans{0}
    \loigiai{Khi $x$ tiến tới vô cùng, $1/x$ tiến tới 0.}
\end{ex}

%Môn: Vật Lí
\begin{ex}
    Cường độ âm chuẩn là $I_0 = 10^{-12} W/m^2$. Nếu cường độ âm tại một điểm là $I = 10^{-10} W/m^2$ thì mức cường độ âm $L$ là bao nhiêu dB?
    \shortans{20}
    \loigiai{$L = 10 \log(I/I_0) = 10 \log(100) = 20$ dB.}
\end{ex}

%Môn: Địa Lí
\begin{ex}
    Năm 2020, quốc gia nào có dân số đông nhất khu vực Mỹ La-tinh?
    \choice
    {Mê-hi-cô}
    {Chi-lê}
    {\True Bra-xin}
    {Ác-hen-ti-na}
    \loigiai{Bra-xin là quốc gia đông dân nhất khu vực này (hơn 211 triệu người năm 2020).}
\end{ex}

%Môn: Sinh Học
\begin{ex}
    Ở thực vật, nước được vận chuyển theo một chiều từ rễ lên lá nhờ:
    \choiceTF
    {\True Lực hút do thoát hơi nước ở lá.}
    {\True Lực đẩy của áp suất rễ.}
    {\True Lực liên kết giữa các phân tử nước.}
    {Mạch rây đóng vai trò vận chuyển chính.}
    \loigiai{
        \begin{itemchoice}
        \itemch đây là động lực chính.
        \itemch 
        \itemch 
        \itemch mạch gỗ mới là nơi vận chuyển nước.
        \end{itemchoice}
    }
\end{ex}

%Môn: Toán
\begin{ex}
    Trong một hình hộp chữ nhật có bao nhiêu cạnh?
    \shortans{12}
    \loigiai{Tương tự hình lập phương, hình hộp chữ nhật có 12 cạnh.}
\end{ex}